
\section{Introduction}

In his 1977 preprint \cite{Ulmer1977Bialgebras}, Ulmer introduced a semantic notion of a bialgebra to unify a broad range of algebraic and coalgebraic structures within the setting of locally presentable categories. Ulmer's insight was to view a bialgebraic structure on an object $c$ as a structural morphism $f\colon Fc\to Gc$ satisfying an equation $\alpha_{(c,f)} = \beta_{(c,f)}$ of natural transformations. Varieties of universal algebra, such as categories of groups or rings, Eilenberg--Moore categories of monads, internal bimonoids or Hopf algebras in symmetric monoidal categories are all of this form. Rather than treating these examples individually, Ulmer showed that many of their fundamental structural properties follow from a common semantic framework. Ulmer left out the syntactic or schematic approach, which he described as ``hopelessly involved'' \cite[3.1]{Ulmer1977Bialgebras}.

Our guiding idea is that the commutative diagrams defining a bialgebraic structure can be understood as syntax for equations between $2$-cells. For example, the commutative diagrams defining Hopf or Frobenius algebras in a symmetric monoidal category can be interpreted as pairs of natural transformations constructed from the data of a symmetric pseudomonoid in a $2$-category with PIE limits. This point of view gives meaning to $T$-bialgebras, for a bialgebraic theory $T$, in $2$-categories whose objects might not have any concrete morphisms of their own.

We develop Ulmer's notion of bialgebras within an abstract $2$-category $\K$ in two different layers. In the first semantic layer, we define a \emph{bialgebraic datum} on an object $c$ in $\K$ as a tuple
$$
\B = (c,\ f,g\colon c\to x,\ f',g'\colon c\to y,\ \alpha,\beta\colon f'u'\Rightarrow g'u')
$$
consisting of two \emph{function support} morphisms $f,g$, whose inserter $u'\colon \Ins(f,g)\to c$ defines Ulmer's \textit{object of pre-bialgebras} with the \textit{equifier} $e\colon \Eq(\alpha,\beta)\to \Ins(f,g)$ defining the object $\BialgB$ of \textit{bialgebras} of $\B$. The structural results of this paper mostly relate to such objects of bialgebras. The second layer is syntactic. A \textit{signature pair} $(\Sigma,\sigma)$ and a bialgebraic $\sigma$-\textit{theory} $T$ describe the structure maps and their equations relative to a $2$-algebraic structure fixed by $\Sigma$, such as a monoidal or a monadic structure. For a fixed $\Sigma$-model $M$ in $\K$, we can construct the bialgebraic datum $\B_{M,T}$, whose object of bialgebras is the object $M^T$ of \textit{internal $T$-bialgebras}.

\smallskip
\noindent\textbf{Main results:} Capital letters $F,G,F',G'$ indicate that $\K = \CAT$.
\begin{itemize}
    \item \textit{Induced Functor of Algebras} (Theorem \ref{thm:induced_morphism_between_algebras}). We define general notions of lax, oplax and strong morphisms $M\to N$ of $\Sigma$-models in $\K$, which induce a $1$-cell $M^T\to N^T$ for algebraic, coalgebraic and bialgebraic theories $T$, respectively. This generalizes the classical lifting of a lax monoidal functor to a functor between the categories of internal monoids.
    
    \item \textit{Strict creation of Kan extensions} (Theorem \ref{thm:pie_kan_creation}, Corollary \ref{cor:strictCreationToBialgebraicDatumObjectsKan}). We show that the canonical morphism $u_\B\colon \BialgB\to c$ strictly creates those left Kan extensions that the arity support $f$ preserves and whose quotient initiality $f'$ preserves. This generalizes the classical result that the forgetful functor $\C^T\to \C$ from the Eilenberg--Moore category of a monad $T$ on $\C$ strictly creates the colimiting cocones that $T$ preserves and whose joint epicness $T^2$ preserves.

    \item \textit{Monadicity and detection of left Kan extensions} (Theorem \ref{thm:BigMonadicityB}, Corollary \ref{cor:BigMonadicityModel}). If the coarity supports $G$ and $G'$ of $\B$ are right adjoints, $\C$ has colimits of size up to an infinite cardinal $\kappa$ and the arity supports $F$ and $F'$ preserve colimits of $\kappa$-chains, then $U\colon \BialgB\to \C$ is monadic and $\BialgB$ has colimits of size at most $\kappa$. These results are then upgraded to the $2$-categorical setting appropriately in Section \ref{sec:representable}.

    \item \textit{Local presentability} (Theorem \ref{thm:LPandComonadicity}). Over a locally presentable base $\C$, $\BialgB$ is locally presentable and $U_\B$ is monadic, given that the support functors are accessible and the coarity support functors $G$ and $G'$ are continuous. The corresponding comonadicity is obtained when the arity functors $F$ and $F'$ are cocontinuous. This theorem is generalized suitably to the $2$-categorical setting in Theorem \ref{thm:2CatLPComonadicity}. The definition of local presentability of an object in a $2$-category is given in Definition \ref{def:ALocalPresentability}.

    \item \textit{Factorization systems, regularity and exactness} (Theorems \ref{thm:pie_ofs_creation} and \ref{thm:liftingregularexact}). If the category $\C$ is equipped with an orthogonal factorization system $(\E,\M)$ such that $F(\E)\perp G(\M)$ and $F'(\E)\perp_! G'(\M)$, then $\BialgB$ has a unique orthogonal factorization system that the forgetful functor $U_\B\colon \BialgB\to \C$ preserves. Combining the limit and orthogonal factorization creation results, we obtain conditions under which regularity and exactness lift from the base category $\C$ to the object $\BialgB$ of bialgebras.
\end{itemize}

\smallskip
\noindent\textbf{Related work:}
Ulmer's work \cite{Ulmer1977Bialgebras} was set in the context of locally presentable categories, where he studied completeness and size properties of bialgebraic structures. Among the examples given were algebras over a PROP and descent data. Ulmer left open the local presentability of arbitrary Hopf algebras over a commutative ring \cite[4.6]{Ulmer1977Bialgebras}. Porst later proved this local presentability \cite{Porst-CoLimits} and studied local presentability of internal monoids in symmetric monoidal settings \cite{Porst01062008} and of Hopf algebras \cite{PorstHopfLocallyPresentable, Porst2}. We also touch on the setting of Hopf algebras by proving the exactness of the category $\Hopf_R^\coc$ of cocommutative Hopf algebras over a commutative von Neumann regular ring $R$ (Corollary \ref{cor:algebraicExtensionLiftingExactness}), extending the result for fields \cite{GRAN20194171}. The same corollary also establishes the exactness of many structures extending cocommutative Hopf algebras, such as $\Omega$-Hopf algebras \cite{bevilacquaCoalgebraicModelsOmega2026} and cocommutative Hopf braces \cite{granHopfBracesSemiabelian2026}. 

Theorem \ref{thm:pie_ofs_creation} shows how to lift orthogonal factorization systems to categories of bialgebras, generalizing results by Linton \cite{Linton1969} and Wi{\ss}mann \cite{Wissmann2022Minimality}. Linton also showed in the same article how to compute colimits in the Eilenberg--Moore category $\C^T$ of a monad, provided that $\C^T$ has reflexive coequalizers. This result was extended by Barr \cite{barrCoequalizersFreeTriples1970} and Ad\'{a}mek \cite{adamekColimitsAlgebrasRevisited1977}. Theorem \ref{thm:BigMonadicityB} further develops these ideas by computing left Kan extensions without completeness and smallness assumptions, but requiring the arity functors' adherence to colimits of $\kappa$-chains.

On a $2$-categorical note, our objects of bialgebras are constructed from PIE limits, which are flexible limits \cite{birdFlexibleLimits2categories1989}. We prove $2$-categorical monadicity results (Theorems \ref{thm:2CategoricalMonadicityForBialgebras}, \ref{thm:2CatLPComonadicity}), which connect to the formal theory of monads \cite{STREET1972149, lackFormalTheoryMonads2002}. The notion of bi-density, which we use to generalize our results to objects of $2$-categories, is a variation of the bicategorical notion of density due to Hofstra and Karvonen \cite{hofstraInnerAutoequivalencesGeneral2024}. The notion of density was introduced by Isbell in \cite{isbellAdequateSubcategories1960}, further developed by Ulmer \cite{ulmerPropertiesDenseRelative1968} and Dubuc \cite{dubuc1970}. 

\smallskip
\noindent\textbf{Outline:} Section~\ref{sec:preliminaries} reviews the required $2$-categorical preliminaries and defines strict products, inserters and equifiers in $2$-categories, so-called PIE limits. Section~\ref{sec:bialgebras} defines bialgebraic data $\B$ and their object $\BialgB$ of bialgebras, and develops signature pairs $(\Sigma,\sigma)$, theories $T$ and $\Sigma$-models $M$, constructs $M^T$, and proves the Induced Functor of Algebras Theorem. Section~\ref{sec:lifting_limits} proves the creation result for Kan extensions and specializes it to sifted colimits and wide pushouts of regular epimorphisms in $\CAT$. Section~\ref{sec:monadicity} studies the detection of colimits and proves monadicity theorems for the forgetful functor $\BialgB\to \C$, together with comonadicity under accessibility assumptions. Section~\ref{sec:ofs} proves the lifting result for orthogonal factorization systems along PIE limits and applies this result to derive corresponding lifting theorems for regularity and exactness. Finally, Section~\ref{sec:representable} uses representability and bi-density to generalize the monadicity and comonadicity results to $2$-categories. The appendix shows that the $\V$-categories $\textbf{1},\textbf{2}_x, \textbf{3}_{x,y}$ for objects $x,y\in \Obj(\V)$ form a bi-dense class of objects in $\V\text{-}\CAT$ when $\V$ is a monoidal category with an absorbing initial object.